% =========================
% L20.tex
% =========================

\chapter{L20 -- Equilibri, movimenti, stabilità e linearizzazione}
\label{chap:L20}

\section{Richiamo iniziale: stato, traiettoria e rappresentazione grafica}

Prima di introdurre i concetti di equilibrio e stabilità, è utile chiarire bene il significato geometrico di una traiettoria di stato.

Consideriamo, come esempio introduttivo, l'oscillatore armonico non forzato
\[
\ddot y(t) + y(t) = 0.
\]
Introducendo le variabili di stato
\[
x_1(t)=y(t), 
\qquad
x_2(t)=\dot y(t),
\]
si ottiene il sistema
\[
\begin{cases}
\dot x_1(t)=x_2(t),\\
\dot x_2(t)=-x_1(t).
\end{cases}
\]
In forma matriciale:
\[
\dot x(t)=
\begin{bmatrix}
0 & 1\\
-1 & 0
\end{bmatrix}x(t).
\]

Se ad esempio si sceglie la condizione iniziale
\[
x(0)=
\begin{bmatrix}
1\\
0
\end{bmatrix},
\]
la soluzione è
\[
x_1(t)=\cos t,
\qquad
x_2(t)=-\sin t.
\]
Quindi, nel piano di stato \((x_1,x_2)\), la traiettoria descrive una circonferenza unitaria.

\subsection*{Interpretazione geometrica}
Questo esempio fa vedere bene la differenza tra:
\begin{itemize}
    \item il \textbf{movimento nel tempo}, cioè l'evoluzione delle componenti \(x_1(t)\), \(x_2(t)\) come funzioni di \(t\);
    \item la \textbf{traiettoria nello spazio di stato}, cioè il luogo dei punti \(x(t)\) nel piano \((x_1,x_2)\).
\end{itemize}

\begin{figure}[h!]
\centering
\begin{tikzpicture}[scale=1.25]
    % assi
    \draw[->] (-2.2,0) -- (2.4,0) node[right] {$x_1$};
    \draw[->] (0,-2.2) -- (0,2.4) node[above] {$x_2$};

    % circonferenza
    \draw[dashed] (0,0) circle (1.5);

    % frecce senso orario
    \draw[->] (1.5,0) arc[start angle=0,end angle=-35,radius=1.5];
    \draw[->] (0,-1.5) arc[start angle=-90,end angle=-125,radius=1.5];
    \draw[->] (-1.5,0) arc[start angle=180,end angle=145,radius=1.5];
    \draw[->] (0,1.5) arc[start angle=90,end angle=55,radius=1.5];

    % punto iniziale
    \fill (1.5,0) circle (1.5pt);
    \node[below right] at (1.5,0) {$x(0)$};

    \node at (1.15,1.35) {\small \(x_1^2+x_2^2=1\)};
\end{tikzpicture}
\caption{Traiettoria dell'oscillatore armonico nel piano di stato.}
\end{figure}

\bigskip

\section{Equilibrio di un sistema dinamico}

Passiamo adesso al concetto di equilibrio.

\subsection{Caso continuo}
Per un sistema non lineare tempo continuo
\[
\dot x = f(x,u),
\]
una coppia \((\bar x,\bar u)\) si dice \textbf{coppia di equilibrio} se, applicando l'ingresso costante \(\bar u\), lo stato \(\bar x\) rimane costante nel tempo. Questo equivale a imporre
\[
f(\bar x,\bar u)=0.
\]

\subsection{Caso discreto}
Per un sistema non lineare tempo discreto
\[
x(k+1)=f(x(k),u(k)),
\]
la coppia \((\bar x,\bar u)\) è di equilibrio se
\[
\bar x=f(\bar x,\bar u).
\]

\subsection*{Interpretazione}
Dire che \((\bar x,\bar u)\) è una coppia di equilibrio significa che, se il sistema parte da \(\bar x\) e l'ingresso viene mantenuto uguale a \(\bar u\), allora lo stato non si muove:
\[
x(t)\equiv \bar x
\quad \text{oppure} \quad
x(k)\equiv \bar x.
\]

In altre parole, la ``velocità'' dello stato è nulla: il sistema resta fermo in quel punto.

\bigskip

\section{Movimento e traiettoria}

Sia assegnata una soluzione del sistema, ad esempio nel continuo
\[
x(t)=\varphi(t,t_0,x_0,u),
\]
oppure nel discreto
\[
x(k)=\varphi(k,k_0,x_0,u).
\]

\begin{itemize}
    \item Il \textbf{movimento} è l'evoluzione temporale del sistema.
    \item La \textbf{traiettoria} è la curva descritta dal vettore di stato nello spazio di stato.
\end{itemize}

Se ad esempio si guarda una sola componente \(x_i(t)\), si ottiene il grafico in funzione del tempo. Se invece si rappresenta \(x(t)\) nello spazio degli stati, si ottiene la traiettoria.

\begin{figure}[h!]
\centering
\begin{tikzpicture}[scale=1.0]
    % sinistra: x(t)
    \draw[->] (0,0) -- (4.4,0) node[right] {$t$};
    \draw[->] (0,0) -- (0,2.8) node[above] {$x(t)$};
    \draw[smooth] plot coordinates {(0.2,0.4) (0.8,1.2) (1.5,1.0) (2.2,0.6) (3.1,1.8) (4.0,2.2)};
    \draw[dashed] (2.2,0) -- (2.2,0.6);
    \node[below] at (2.2,0) {$t^\star$};

    % destra: traiettoria nello spazio di stato
    \begin{scope}[xshift=7cm]
        \draw[->] (-2.2,0) -- (2.5,0) node[right] {$x_1$};
        \draw[->] (0,-2.0) -- (0,2.5) node[above] {$x_2$};
        \draw[smooth] plot coordinates {(-1.8,-1.2) (-1.0,-0.5) (-0.3,0.8) (0.4,1.3) (1.0,0.7) (1.8,1.8)};
        \fill (0.4,1.3) circle (1.2pt);
        \node[right] at (0.45,1.3) {\small \(x(t^\star)\)};
    \end{scope}
\end{tikzpicture}
\caption{A sinistra un movimento nel tempo, a destra la corrispondente traiettoria nello spazio di stato.}
\end{figure}

\bigskip

\section{Stabilità di un movimento nominale}

La teoria della stabilità non si applica solo ai punti di equilibrio: in generale si può parlare anche di stabilità di un \textbf{movimento nominale}.

Supponiamo di avere una traiettoria nominale
\[
x(t)=\varphi(t,t_0,x_0,\bar u)
\]
generata da un certo ingresso assegnato \(\bar u\), e una traiettoria perturbata
\[
\tilde x(t)=\varphi(t,t_0,x_0+\delta x_0,\bar u)
\]
ottenuta modificando leggermente la condizione iniziale.

\subsection{Movimento stabile}
Il movimento nominale si dice \textbf{stabile} se
\[
\forall \varepsilon>0 \ \exists \delta>0 \ \text{tale che}\
\|\delta x_0\|<\delta
\ \Longrightarrow\
\|\tilde x(t)-x(t)\|<\varepsilon
\quad \forall t\ge t_0.
\]

Interpretazione: se parto sufficientemente vicino alla traiettoria nominale, allora ci resto vicino per tutti i tempi successivi.

\subsection{Movimento attrattivo}
Il movimento nominale si dice \textbf{attrattivo} se esiste un intorno iniziale tale che
\[
\|\delta x_0\|<\delta
\ \Longrightarrow\
\lim_{t\to+\infty}\|\tilde x(t)-x(t)\|=0.
\]

Interpretazione: le traiettorie perturbate, magari dopo essersi un po' allontanate, alla fine tendono alla traiettoria nominale.

\subsection{Movimento asintoticamente stabile}
Il movimento nominale si dice \textbf{asintoticamente stabile} se è sia stabile sia attrattivo.

\bigskip

\section{Esempio: attrattivo ma non stabile}

Un esempio classico, presente anche nei tuoi appunti, è il sistema discreto scalare
\[
x(k+1)=
\begin{cases}
2x(k), & |x(k)|<1,\\[0.1cm]
0, & |x(k)|\ge 1.
\end{cases}
\]

Studiamo l'equilibrio nell'origine \(x=0\).

\subsection*{Perché è attrattivo}
Se la traiettoria prima o poi esce dall'intervallo \((-1,1)\), al passo successivo va a zero e poi resta a zero. Dunque molte traiettorie convergono all'origine.

\subsection*{Perché non è stabile}
Prendiamo una condizione iniziale molto piccola, ad esempio
\[
x(0)=\frac{1}{2}.
\]
Allora
\[
x(1)=2x(0)=1,
\qquad
x(2)=0.
\]
Quindi la traiettoria, pur convergendo infine all'origine, prima si allontana da un intorno arbitrariamente piccolo dello zero. Questo viola la definizione di stabilità.

\begin{figure}[h!]
\centering
\begin{tikzpicture}[scale=1.0]
    \draw[->] (0,0) -- (5.2,0) node[right] {$k$};
    \draw[->] (0,-0.4) -- (0,2.6) node[above] {$x(k)$};

    % livelli
    \draw[dashed] (0,1) -- (4.7,1);
    \draw[dashed] (0,0.5) -- (4.7,0.5);

    % punti
    \fill (0,0.5) circle (1.6pt);
    \fill (1,1.0) circle (1.6pt);
    \fill (2,0.0) circle (1.6pt);
    \fill (3,0.0) circle (1.6pt);
    \fill (4,0.0) circle (1.6pt);

    \draw (0,0.5) -- (1,1.0) -- (2,0.0) -- (4,0.0);

    \node[below] at (0,0) {$0$};
    \node[below] at (1,0) {$1$};
    \node[below] at (2,0) {$2$};
\end{tikzpicture}
\caption{Traiettoria che converge a zero ma non resta vicina allo zero per tutti i tempi: l'origine è attrattiva ma non stabile.}
\end{figure}

\bigskip

\section{Equilibrio come caso particolare di movimento}

Un punto di equilibrio è un caso particolare di traiettoria nominale: è una traiettoria costante.

Infatti, se \((\bar x,\bar u)\) è una coppia di equilibrio, allora la traiettoria nominale associata è
\[
x(t)\equiv \bar x
\qquad \text{oppure} \qquad
x(k)\equiv \bar x.
\]

Per questo motivo le definizioni di stabilità, attrattività e stabilità asintotica di un equilibrio si ottengono come caso particolare di quelle date per un movimento qualunque.

\bigskip

\section{Traslazione dell'equilibrio all'origine}

Quando si studia la stabilità di un equilibrio \(\bar x\), è quasi sempre conveniente traslare il sistema in nuove coordinate in modo che l'equilibrio diventi l'origine.

Definiamo
\[
z=x-\bar x,
\qquad
v=u-\bar u.
\]

\subsection{Caso continuo}
Se
\[
\dot x=f(x,u),
\]
allora
\[
\dot z=\dot x=f(z+\bar x,v+\bar u).
\]
Poiché \((\bar x,\bar u)\) è equilibrio, vale
\[
f(\bar x,\bar u)=0.
\]
Quindi il sistema traslato ha equilibrio nell'origine \(z=0\).

\subsection{Caso discreto}
Se
\[
x(k+1)=f(x(k),u(k)),
\]
ponendo
\[
z(k)=x(k)-\bar x,
\qquad
v(k)=u(k)-\bar u,
\]
si ottiene di nuovo un sistema equivalente con equilibrio nell'origine.

\medskip

\noindent
\textbf{Conclusione importante.}
Le proprietà di stabilità dell'equilibrio \(\bar x\) nel sistema originale coincidono con le proprietà di stabilità dell'origine nel sistema traslato.

Per questo, da qui in poi, si studia quasi sempre la stabilità dell'origine.

\bigskip

\section{Equilibri dei sistemi lineari autonomi}

Consideriamo ora un sistema lineare autonomo.

\subsection{Tempo continuo}
\[
\dot x=Ax.
\]
Un punto \(\bar x\) è equilibrio se
\[
A\bar x=0.
\]
Quindi l'insieme degli equilibri è
\[
\ker A.
\]

\subsection{Tempo discreto}
\[
x(k+1)=Ax(k).
\]
Un punto \(\bar x\) è equilibrio se
\[
\bar x=A\bar x
\quad \Longleftrightarrow \quad
(A-I)\bar x=0.
\]
Quindi l'insieme degli equilibri è
\[
\ker(A-I).
\]

\subsection*{Osservazione}
Se \(A\) non è invertibile (nel continuo), oppure se \(1\) è autovalore di \(A\) (nel discreto), allora possono esistere infiniti equilibri.

\bigskip

\section{Perché un equilibrio non isolato non può essere asintoticamente stabile}

Se un equilibrio non è isolato, allora in ogni suo intorno esistono altri punti di equilibrio.

Ma una traiettoria che parte esattamente da un altro equilibrio resta ferma in quel punto e quindi non può convergere all'equilibrio considerato. Perciò un equilibrio non isolato non può essere attrattivo, e dunque non può essere asintoticamente stabile.

\medskip

\noindent
Può però essere \textbf{stabile} nel senso di Lyapunov.

\bigskip

\section{Regione di attrazione o RAS}

Sia l'origine un equilibrio asintoticamente stabile. Si definisce \textbf{regione di attrazione} (RAS, \emph{Region of Asymptotic Stability}) l'insieme delle condizioni iniziali che generano traiettorie convergenti all'origine:
\[
\mathcal{A}(0)=\left\{x_0 \in \mathbb{R}^n \ :\ \lim_{t\to +\infty}\varphi(t,t_0,x_0)=0 \right\}.
\]

Se l'origine è \textbf{globalmente asintoticamente stabile}, allora
\[
\mathcal{A}(0)=\mathbb{R}^n.
\]

In generale, però, la RAS può essere un sottoinsieme proprio dello spazio di stato e non è sempre facile da determinare esplicitamente.

\bigskip

\section{Linearizzazione di un sistema non lineare}

Per studiare la stabilità locale di un equilibrio di un sistema non lineare, si usa la linearizzazione.

Consideriamo il sistema
\[
\dot x = f(x,u),
\qquad
y=h(x,u),
\]
e una coppia di equilibrio \((\bar x,\bar u)\), cioè
\[
f(\bar x,\bar u)=0.
\]

Introduciamo le variabili di deviazione:
\[
z=x-\bar x,
\qquad
v=u-\bar u,
\qquad
\eta=y-\bar y,
\quad \bar y=h(\bar x,\bar u).
\]

Sviluppando \(f\) e \(h\) al primo ordine attorno a \((\bar x,\bar u)\) si ottiene
\[
f(x,u)\approx
f(\bar x,\bar u)
+\left.\frac{\partial f}{\partial x}\right|_{(\bar x,\bar u)}(x-\bar x)
+\left.\frac{\partial f}{\partial u}\right|_{(\bar x,\bar u)}(u-\bar u),
\]
e analogamente
\[
h(x,u)\approx
h(\bar x,\bar u)
+\left.\frac{\partial h}{\partial x}\right|_{(\bar x,\bar u)}(x-\bar x)
+\left.\frac{\partial h}{\partial u}\right|_{(\bar x,\bar u)}(u-\bar u).
\]

Poiché \(f(\bar x,\bar u)=0\), il sistema linearizzato diventa
\[
\begin{cases}
\dot z = Az + Bv,\\
\eta = Cz + Dv,
\end{cases}
\]
dove
\[
A=\left.\frac{\partial f}{\partial x}\right|_{(\bar x,\bar u)},
\qquad
B=\left.\frac{\partial f}{\partial u}\right|_{(\bar x,\bar u)},
\qquad
C=\left.\frac{\partial h}{\partial x}\right|_{(\bar x,\bar u)},
\qquad
D=\left.\frac{\partial h}{\partial u}\right|_{(\bar x,\bar u)}.
\]

\medskip

\noindent
\textbf{Interpretazione.}
La linearizzazione fornisce una buona approssimazione del sistema non lineare \emph{solo localmente}, cioè per stati e ingressi sufficientemente vicini al punto di equilibrio considerato.

\bigskip

\section{Esempio: pendolo con attrito}

Consideriamo il pendolo con attrito viscoso e ingresso di coppia:
\[
m\ell^2\ddot\theta + c\dot\theta + mg\ell \sin\theta = u.
\]

Introduciamo le variabili di stato
\[
x_1=\theta,
\qquad
x_2=\dot\theta.
\]
Allora il sistema si scrive come
\[
\begin{cases}
\dot x_1 = x_2,\\
\dot x_2 = -\dfrac{g}{\ell}\sin x_1 - \dfrac{c}{m\ell^2}x_2 + \dfrac{1}{m\ell^2}u.
\end{cases}
\]

\subsection{Equilibri per \(u=0\)}
Se \(u=0\), gli equilibri soddisfano
\[
\begin{cases}
0=x_2,\\
0=-\dfrac{g}{\ell}\sin x_1.
\end{cases}
\]
Quindi
\[
x_2=0,
\qquad
\sin x_1=0
\quad\Longrightarrow\quad
x_1=k\pi,\qquad k\in\mathbb{Z}.
\]
Dunque ci sono infiniti equilibri:
\[
(k\pi,0), \qquad k\in\mathbb{Z}.
\]

\subsection{Jacobiana del sistema}
La matrice Jacobiana rispetto allo stato è
\[
\frac{\partial f}{\partial x}(x_1,x_2)=
\begin{bmatrix}
0 & 1\\[0.1cm]
-\dfrac{g}{\ell}\cos x_1 & -\dfrac{c}{m\ell^2}
\end{bmatrix},
\]
mentre rispetto all'ingresso
\[
\frac{\partial f}{\partial u}(x_1,x_2)=
\begin{bmatrix}
0\\[0.1cm]
\dfrac{1}{m\ell^2}
\end{bmatrix}.
\]

\subsection{Linearizzazione attorno a \((0,0)\)}
Poiché \(\cos 0=1\), si ottiene
\[
A_1=
\begin{bmatrix}
0 & 1\\[0.1cm]
-\dfrac{g}{\ell} & -\dfrac{c}{m\ell^2}
\end{bmatrix},
\qquad
B_1=
\begin{bmatrix}
0\\[0.1cm]
\dfrac{1}{m\ell^2}
\end{bmatrix}.
\]

La traccia e il determinante valgono
\[
\operatorname{tr}(A_1)=-\frac{c}{m\ell^2}<0,
\qquad
\det(A_1)=\frac{g}{\ell}>0.
\]
Quindi gli autovalori hanno parte reale negativa: l'equilibrio \((0,0)\) risulta localmente asintoticamente stabile se \(c>0\).

\subsection{Linearizzazione attorno a \((\pi,0)\)}
Poiché \(\cos \pi=-1\), si ottiene
\[
A_2=
\begin{bmatrix}
0 & 1\\[0.1cm]
\dfrac{g}{\ell} & -\dfrac{c}{m\ell^2}
\end{bmatrix},
\qquad
B_2=
\begin{bmatrix}
0\\[0.1cm]
\dfrac{1}{m\ell^2}
\end{bmatrix}.
\]

Ora
\[
\det(A_2)=-\frac{g}{\ell}<0.
\]
Quindi gli autovalori hanno segno opposto e l'equilibrio \((\pi,0)\) è instabile.

\medskip

\noindent
Poiché il pendolo è periodico in \(\theta\), gli equilibri con \(x_1=2k\pi\) hanno lo stesso comportamento di \((0,0)\), mentre quelli con \(x_1=(2k+1)\pi\) hanno lo stesso comportamento di \((\pi,0)\).

\bigskip

\section{Teorema indiretto di Lyapunov}

Il teorema indiretto di Lyapunov consente di dedurre la stabilità locale del sistema non lineare a partire dalla stabilità del sistema linearizzato.

Consideriamo il sistema autonomo
\[
\dot x=f(x),
\qquad
f(0)=0.
\]
Sia
\[
\dot z = Az
\]
la linearizzazione di \(\dot x=f(x)\) attorno all'origine, con
\[
A=\left.\frac{\partial f}{\partial x}\right|_{x=0}.
\]

Allora valgono i seguenti risultati.

\subsection*{1. Caso asintoticamente stabile}
Se la matrice \(A\) è di Hurwitz, cioè tutti i suoi autovalori hanno parte reale strettamente negativa, allora l'origine del sistema non lineare è \textbf{localmente asintoticamente stabile}.

\subsection*{2. Caso instabile}
Se \(A\) possiede almeno un autovalore con parte reale positiva, allora l'origine del sistema non lineare è \textbf{instabile}.

\subsection*{3. Caso critico}
Se tutti gli autovalori di \(A\) hanno parte reale non positiva, ma almeno uno ha parte reale nulla, allora il teorema indiretto \textbf{non è conclusivo}. In questo caso servono altri strumenti, per esempio il metodo diretto di Lyapunov.

\bigskip

\section{Caso critico: perché il teorema indiretto può non bastare}

Se, nel pendolo, si pone \(c=0\), la linearizzazione attorno a \((0,0)\) diventa
\[
A=
\begin{bmatrix}
0 & 1\\
-\dfrac{g}{\ell} & 0
\end{bmatrix}.
\]
Allora
\[
\operatorname{tr}(A)=0,
\qquad
\det(A)=\frac{g}{\ell}>0,
\]
e gli autovalori sono
\[
\lambda_{1,2}=\pm j\sqrt{\frac{g}{\ell}}.
\]

Il sistema linearizzato ha quindi orbite periodiche chiuse, e il teorema indiretto di Lyapunov non permette di concludere né stabilità asintotica né instabilità del sistema non lineare.

Questo è il tipico \textbf{caso critico}.

\bigskip

\section{Schema riassuntivo sulle classificazioni di stabilità}

Per un equilibrio dell'origine nel continuo:
\begin{itemize}
    \item se esiste un autovalore con \(\Re(\lambda)>0\), allora l'origine è \textbf{instabile};
    \item se tutti gli autovalori hanno \(\Re(\lambda)<0\), allora l'origine è \textbf{localmente asintoticamente stabile};
    \item se tutti hanno \(\Re(\lambda)\le 0\) e qualcuno ha \(\Re(\lambda)=0\), il test lineare è \textbf{inconcludente}.
\end{itemize}

Nel discreto, le condizioni si riscrivono sostituendo:
\[
\Re(\lambda)<0
\quad \longleftrightarrow \quad
|\lambda|<1,
\]
\[
\Re(\lambda)>0
\quad \text{(nel caso continuo)}
\quad \leadsto \quad
|\lambda|>1
\quad \text{(nel caso discreto)}.
\]

\bigskip

\section{Esempio geometrico: equilibrio stabile e instabile}

Un modo utile per interpretare la stabilità è osservare il comportamento delle traiettorie vicine a un punto di equilibrio.

\begin{itemize}
    \item Se tutte le traiettorie che partono abbastanza vicine restano vicine, l'equilibrio è \textbf{stabile}.
    \item Se inoltre tendono all'equilibrio, è \textbf{asintoticamente stabile}.
    \item Se esistono traiettorie arbitrariamente vicine che se ne allontanano, l'equilibrio è \textbf{instabile}.
\end{itemize}

Ad esempio, nel pendolo:
\begin{itemize}
    \item la posizione verso il basso \(\theta=0\) è stabile (e con attrito è anche asintoticamente stabile);
    \item la posizione capovolta \(\theta=\pi\) è instabile.
\end{itemize}

\bigskip

\section{Conclusioni della lezione}

In questa lezione abbiamo chiarito alcuni concetti fondamentali per lo studio qualitativo dei sistemi dinamici.

\begin{enumerate}
    \item Il concetto di \textbf{equilibrio} è il caso particolare di una traiettoria costante.
    \item La stabilità può essere definita inizialmente per un \textbf{movimento nominale} e poi specializzata al caso dell'equilibrio.
    \item Un equilibrio si studia più comodamente traslandolo nell'origine.
    \item Nei sistemi lineari autonomi, l'insieme degli equilibri è
    \[
    \ker A \quad \text{(TC)},
    \qquad
    \ker(A-I) \quad \text{(TD)}.
    \]
    \item Un equilibrio non isolato non può essere asintoticamente stabile.
    \item La \textbf{RAS} è l'insieme delle condizioni iniziali che convergono all'equilibrio.
    \item Per i sistemi non lineari, la \textbf{linearizzazione} permette di dedurre proprietà locali di stabilità.
    \item Il \textbf{teorema indiretto di Lyapunov} consente di concludere:
    \begin{itemize}
        \item stabilità asintotica locale se il linearizzato è asintoticamente stabile;
        \item instabilità se il linearizzato ha almeno un modo divergente;
        \item nessuna conclusione nel caso critico.
    \end{itemize}
\end{enumerate}

\medskip

In sostanza, la linearizzazione è il ponte che collega lo studio locale dei sistemi non lineari alla teoria, molto più semplice e potente, dei sistemi lineari.